\section{B4: attainable energy closure and a valid compactness criterion}
\label{sec:r33-b4}

The fast-collision counterexample defeats a state that retains the direction of recession energy while claiming uniform continuous dynamics under a second-moment bound. Collapsing that direction is not, by itself, a proof of a boundary generator. This chapter characterizes a smaller kinematic closure and proves a positive path-compactness theorem under explicit tail control. It keeps the original running action and does not insert a penalty into it.

\subsection{The actual narrow--total-energy closure}
Let $S=\mathbb T^d\times\mathbb R^d$ and $E(f)=\int|v|^2df$. In the product of narrow probability convergence and ordinary real convergence define
\[
 \mathcal K_{E_0}=\{(f,e): f\in\mathcal P(S),\quad E(f)\le e\le E_0\}.
\]

\begin{proposition}[Attainability and absence of mass at infinity]
\label{prop:r33-b4-closure}
$\mathcal K_{E_0}$ is compact and is exactly the closure of
$\{(f,E(f)):E(f)\le E_0\}$. The defect $e-E(f)$ can be positive, but positive probability mass at infinite velocity is not attainable under the same bound. Momentum converges automatically along these sequences.
\end{proposition}

\begin{proof}
Markov's inequality gives $f(|v|>R)\le E_0/R^2$, so the probabilities are tight. The interval for $e$ is compact and lower semicontinuity of $E$ makes the displayed set closed. For attainability take $(f,e)$, put $m=E(f)$, and select a probability $g_R$ with symmetric velocities of size $R$ and energy $E(g_R)=R^2$. If $e>m$ choose $p_R=(e-m)/(R^2-m)$ and $f_R=(1-p_R)f+p_Rg_R$. Then $E(f_R)=e$ and $f_R\Rightarrow f$. For $e=m$ use the constant sequence. Smooth tail measures may replace the two atoms by adjusting their energy in the same formula. Finally $\int_{|v|>R}|v|df\le E_0/R$ proves uniform integrability of momentum and its convergence.
\end{proof}

This space records only total escaped energy, not its rapidly changing position or direction. The proposition is kinematic. It defines no evolution for the missing energy and hence is not labelled a semigroup theorem.

\subsection{A path theorem with the missing hypotheses supplied explicitly}
Fix a nonnegative continuous increasing superlinear function $\Psi$, meaning $\Psi(r)/r\to\infty$. Consider density paths $f_t$ satisfying the weak kinetic balance with a nonnegative contact measure $\Gamma=qA_f$, where the hard-sphere angular factor is integrable and bounded by a constant times $|v-v_*|$. Assume uniformly over the family
\[
 \sup_t\int\Psi(|v|^2)f_t\,dx\,dv\le L_\Psi,\qquad
 \sup_{t,x}\rho_t(x)\le R_0,\qquad E(f_t)\le E_0,
\]
where $\rho_t(x)=\int f_t(x,v)dv$, and assume the original running cost
\[
 J_T=\int\ell(q)dA_f\le L.
\]
These are independent hypotheses. A Gaussian entropy bound does not prove the first one.

\begin{theorem}[Entropy-controlled narrow modulus and $W_2$ compactness]
\label{thm:r33-b4-compact}
This family is relatively compact in $C([0,T];\mathcal P_2(S))$ with the $W_2$ topology, provided the weak-balance representatives are used. In particular all limiting density states have zero energy defect.
\end{theorem}

\begin{proof}
The local density bound and $|v-v_*|\le|v|+|v_*|$ give
$A_f([s,t]\times\text{contacts})\le C_A|t-s|$, where $C_A$ depends only on the displayed constants. The entropy inequality $\lambda q\le\ell(q)+e^\lambda-1$ gives for any interval $I$ of length $h$,
\[
 \Gamma(I)\le\frac{L+C_Ah(e^\lambda-1)}\lambda.
\]
Choose $\lambda=\frac12\log(1/h)$ for sufficiently small $h$. The right side tends to zero uniformly. A mass test $\phi$ with bounded-Lipschitz norm at most one changes under free transport by at most $h\sqrt{E_0}$ and under collisions by at most $4\Gamma(I)$. The weak balance therefore yields a uniform bounded-Lipschitz time modulus.

The $\Psi$ bound gives uniform integrability of $|v|^2$, since the tail is bounded by $L_\Psi\sup_{r\ge R^2}r/\Psi(r)$. At each time the family is thus relatively compact in $W_2$. On its compact $W_2$ closure, the narrow and $W_2$ topologies coincide and the inverse identity is uniformly continuous. The narrow modulus consequently gives a $W_2$ modulus. Metric Arzela--Ascoli proves the assertion.
\end{proof}

The counterexample $f_R=(1-R^{-2})M+R^{-2}M(\cdot-Re_1)$ has bounded Gaussian relative entropy and bounded energy, but violates the uniform $\Psi$ bound for every superlinear $\Psi$. The theorem does not exclude it by silently redefining the original entropy action; it states exactly the extra estimate a dynamic application must prove.

\subsection{Collision products need spatial strong convergence}
Even $W_2$ convergence is not enough for the local product $f(x,v)f(x,v_*)$. On a spatial torus take
$f_n(x,v)=(1+a\sin(2\pi nx_1))g(v)$ with $0<a<1$ and a fixed smooth compactly supported velocity density $g$. These probabilities converge to $g(v)dx$ in $W_2$. But the spatially integrated collision product has factor
$\int(1+a\sin(2\pi nx_1))^2dx=1+a^2/2$, not one.

A sufficient positive replacement is strong weighted $L^1$ convergence together with
$\sup_{n,x}\int(1+|v|)f_n(x,v)dv\le R_1$. In fact
\[
 \int (1+|v|)(1+|v_*|)|f_nf_{n,*}-ff_*|\,dx\,dv\,dv_*
 \le2R_1\|f_n-f\|_{L^1(1+|v|)}.
\]
Expand the difference into $(f_n-f)f_{n,*}+f(f_{n,*}-f_*)$ and integrate to prove this inequality. It controls the hard-sphere collision factor and proves collision-product convergence under genuinely sufficient assumptions.

\subsection{Preparation is charged once; semigroup claims need dynamics}
For a fixed initial state $x$ let
\[
 J_t^x(f,\Gamma)=\int_0^t\ell(q)dA_f,\quad f_0=x,
 \qquad V_t\phi(x)=\sup_{f_0=x}\{\phi(f_t)-J_t^x(f,\Gamma)\}.
\]
If the admissible path family is nonempty, stable under cutting and concatenation, and measurable selection of arbitrarily near optimal continuations is available, then dynamic programming gives $V_0=I$ and $V_{t+s}=V_tV_s$. Both inequalities follow by cutting a competitor or concatenating a near optimizer. Preparation entropy appears separately as $\sup_x[V_t\phi(x)-H(x\mid f_0^{\rm ref})]$.

For bounded $\phi$, only controls of cost at most $2\|\phi\|_\infty+1$ need be considered when a zero-cost path exists. A uniform small-time modulus for those controls on a specified invariant state set then proves uniform strong continuity there. Neither state-set invariance nor this modulus is a consequence of the algebraic semigroup identity. On $\mathcal K_{E_0}$ a boundary evolution and its contact cost would still have to be defined and proved compatible with microscopic limits.

\paragraph{Mechanical application still to be established.}
The programme retains its full finite-action kinetic semigroup objective. The outstanding alternatives are to prove genuine uniform tail and spatial compactness estimates for the original controls, or to construct and identify a boundary collision dynamics including its action. No viscosity comparison, arbitrary-boundary realization or full microscopic semigroup limit is claimed from kinematic compactness alone.
